% ============================================================================
% Appendix: Error Calculation in Plots
%
% Required packages (add to preamble if not already present):
%   \usepackage{amsmath}
%   \usepackage{amssymb}
%   \usepackage{verbatim}     % for the ASCII flow diagram in \begin{verbatim}
% ============================================================================

\chapter{Error Calculation in Plots}
\label{app:error-calculation}

All shaded bands (and error bars) shown in the averaged compliance, displacement,
and strain plots represent the \emph{standard error of the mean} (SEM) computed
\emph{across replicate test runs}, not within a single run. The full procedure
is implemented in \texttt{src/data\_utils.py} (averaging) and
\texttt{src/plotting\_functions.py} (rendering).

\section{Aggregating Replicate Runs onto a Common $x$-grid}
\label{sec:err-align}

For a given folder of replicate tests, each run produces a time series
$(x_i, y_i)$, where $x_i$ is the independent variable (e.g.\ test time) and
$y_i$ is the dependent variable (e.g.\ creep compliance). Because runs are
sampled at slightly different time points, the data must be aligned before any
cross-run statistic can be computed. Two alignment modes are supported.

\paragraph{Mode 1 -- Binned averaging}
(\texttt{average\_data\_with\_bins}, \texttt{sem\_by\_run=True}):
\begin{itemize}
  \item Construct $B$ bin edges across the pooled $x$-range
        (linear by default, log-spaced if \texttt{--log\_x} is set, or
        quantile-based if \texttt{--adaptive} is set).
  \item For each run independently, group its samples by bin and take the
        per-bin mean. This yields a matrix
        $Y \in \mathbb{R}^{R \times B}$, where $R$ is the number of runs and
        $B$ is the number of bins.
\end{itemize}

\paragraph{Mode 2 -- Interpolated averaging}
(\texttt{average\_data\_no\_bins}):
\begin{itemize}
  \item Pick the longest valid run as the reference $x$-grid $x_{\text{ref}}$,
        after excluding any run whose $x$-span is less than $50\%$ of the
        median span (truncated runs are dropped).
  \item Linearly interpolate each remaining run onto $x_{\text{ref}}$
        (NaN outside its native range).
  \item This again yields a matrix $Y \in \mathbb{R}^{R \times B}$ with
        $B = |x_{\text{ref}}|$.
\end{itemize}

In both modes, each column of $Y$ contains the values of all runs at one
$x$-location.

\section{Per-$x$ Mean, Standard Deviation, and Sample Size}
\label{sec:err-perx}

For each column $j$ (each bin or each grid point) we compute
%
\begin{equation}
  \bar{y}_j \;=\; \frac{1}{n_j} \sum_{i \in V_j} Y_{ij},
  \qquad
  s_j \;=\; \sqrt{\frac{1}{n_j} \sum_{i \in V_j}
                  \bigl(Y_{ij} - \bar{y}_j\bigr)^2},
  \qquad
  n_j \;=\; |V_j|,
  \label{eq:per-x-stats}
\end{equation}
%
where $V_j$ is the set of runs that contribute a finite value at column $j$.
The standard deviation uses the population convention
($\text{ddof} = 0$, i.e.\ divide by $n_j$ rather than $n_j - 1$); see
\texttt{np.nanstd(..., ddof=0)} in \texttt{average\_data\_with\_bins} and
\texttt{average\_data\_no\_bins}. These three quantities are stored as
\verb|<col>_mean|, \verb|<col>_std|, and \verb|<col>_n| on the returned
DataFrame.

\section{Conversion from Standard Deviation to Standard Error}
\label{sec:err-sem}

The plotted band is the standard error of the mean, obtained from the per-column
standard deviation by dividing by the square root of the run count. This is
done in \texttt{\_std\_to\_sem} in \texttt{src/plotting\_functions.py}:
%
\begin{equation}
  \mathrm{SEM}_j \;=\; \frac{s_j}{\sqrt{n_j}}.
  \label{eq:sem}
\end{equation}
%
When $n_j < 2$ the SEM is undefined and is set to NaN (no band is drawn at that
$x$). The SEM characterises the uncertainty of the \emph{across-run mean},
i.e.\ how reproducibly the cohort estimates the central tendency of the
underlying response. This is distinct from the spread of individual runs
(the standard deviation). SEM is the appropriate metric when the quantity of
interest is the mean compliance/displacement curve of the material rather than
the run-to-run variability itself.

\section{What Is Rendered}
\label{sec:err-render}

For each averaged trace the plot shows:
%
\begin{itemize}
  \item the \textbf{line} at $\bar{y}_j$ versus $x_j$ (per-$x$ cohort mean);
  \item a \textbf{shaded band} between
        $\bar{y}_j - k \cdot \mathrm{SEM}_j$ and
        $\bar{y}_j + k \cdot \mathrm{SEM}_j$, where $k$ is the optional
        \texttt{--error\_scale} multiplier
        (default $k = 1$, i.e.\ $\pm 1\,\mathrm{SEM}$, an approximate $68\%$
        confidence band under approximate normality of the cohort mean).
        In free-style scatter mode (\texttt{--fs} / \texttt{--fs\_lines}) the
        band is replaced by matching error bars.
\end{itemize}

Passing \texttt{--nosem} swaps the SEM band for the raw standard deviation
$s_j$, and \texttt{--error\_scale K} scales the band by an arbitrary factor
$K$ (e.g.\ $K = 1.96$ for an approximate $95\%$ interval).

\section{Summary}
\label{sec:err-summary}

\begin{verbatim}
Replicate runs --> align (bin or interpolate) --> Y in R^{R x B}
                                                    |
                                ,-------------------+-------------------.
                                v                   v                   v
                         mean y_j            std s_j (ddof=0)        n_j (>= 2)
                                                    |
                                                    v
                                          SEM_j = s_j / sqrt(n_j)
                                                    |
                                                    v
                                       plot: y_j +/- k * SEM_j  (k = error_scale)
\end{verbatim}

The convention applied to every averaged figure in this work is therefore:
lines are cohort means across replicates, and shaded bands are
$\pm 1\,\mathrm{SEM}$ of that cohort mean, computed independently at each
$x$-location.
